\documentclass[11pt]{article}

\usepackage[margin=1in]{geometry}
\usepackage{amsmath, amssymb}
\usepackage{hyperref}
\usepackage{enumitem}

\title{A Formally Grounded Accounting Kernel for High-Performance Financial Systems}
\author{}
\date{December 2025}

\begin{document}
\maketitle

\begin{abstract}
Modern financial accounting systems must simultaneously deliver high performance,
strict correctness guarantees, auditability, reproducibility, and adaptability to
evolving business rules. In practice, these requirements are often addressed through
large, expressive application codebases that entangle domain semantics with core
accounting mechanics, complicating both correctness reasoning and performance
optimization.

We propose a formally grounded architecture in which economic computation is expressed
as conservation over integer-valued vector spaces, enforced by a minimal, deterministic
accounting kernel. All business semantics are externalized into a versioned,
content-addressed control plane compiled into immutable artifacts. This separation
enables high assurance, deterministic replay, and auditability without sacrificing
throughput or extensibility.
\end{abstract}

\section{Motivation and Problem Statement}

Financial accounting systems are foundational infrastructure for modern financial
institutions. They are expected to provide absolute correctness while operating at
large scale and under continuously evolving regulatory and business constraints.
However, most production systems interleave accounting invariants with business logic,
procedural workflows, and mutable configuration, making correctness difficult to
reason about and performance unpredictable.

This work proposes an alternative architectural approach: a small, non-expressive
accounting kernel that enforces invariants mechanically, paired with a declarative,
versioned control plane that defines all higher-level semantics. The central design
principle is that economic computation should be treated as a system of conserved
quantities rather than as ad hoc procedural logic.

\section{Accounting as Conservation Laws}

At the core of the system, accounting state is modeled as integer-valued vectors over a
dynamically extensible basis of dimensions. Each dimension represents a conserved
quantity, such as currency units, positions, tax lots, or profit-and-loss components.
Transactions are expressed as balanced vectors over this space, generalizing classical
double-entry accounting to arbitrary conserved dimensions.

Ledger state is derived by deterministically folding an append-only transaction log
using associative (monoidal) aggregation. Because all valid transactions conserve
value across the relevant dimensions, the system provides intrinsic guarantees of
internal consistency, replayability, and auditability independent of business logic.

\section{A Minimal, Deterministic Kernel}

The accounting kernel is intentionally non-expressive. Its sole responsibility is to
enforce mechanical invariants such as balance, conservation, and deterministic
aggregation. It encodes no business semantics and exposes no mutable configuration.

To maximize throughput and predictability, the kernel is implemented using
cache-aligned, fixed-width binary records and branch-minimal aggregation loops.
Performance is therefore dominated by memory bandwidth rather than control flow
complexity. Floating-point arithmetic is avoided entirely in favor of fixed-precision
integer representations, ensuring exact conservation and deterministic replay across
hardware architectures.

\section{Externalized Semantics via a Control Plane}

All domain semantics---including posting rules, account mappings, interpretation
policies, and workflow composition---are defined outside the kernel in a declarative
control plane. Control-plane specifications are treated as immutable computational
inputs rather than mutable configuration.

Each control-plane artifact is authored in a human-oriented declarative form, compiled
into a canonical typed representation, and published as a content-addressed,
cryptographically verifiable artifact. Artifacts are frozen per accounting epoch, such
as a net asset value (NAV) period. Every transaction produced by the system explicitly
references the hash of the control-plane artifact under which it was generated.

This design eliminates semantic race conditions and enables complete provenance
reconstruction. Any ledger state can be replayed exactly as computed, under the precise
semantic rules in force at the time of execution.

\section{External Data as Typed Facts}

External data dependencies, including prices, foreign exchange rates, and reference
data, are modeled as first-class typed facts with explicit provenance, temporal
semantics, and source policies. Data acquisition is not implicit; instead, it is
expressed as an explicit effect plan describing a dependency graph of I/O operations.

Effect plans are annotated with latency, cost, and concurrency constraints, allowing
performance and scalability to be reasoned about declaratively. Determinism is preserved
by binding all computations to immutable fact sets that are explicitly referenced by
ledger entries.

\section{Role of Large Language Models}

Large language models are integrated strictly as assistive tools. Their role is limited
to translating human intent into declarative specifications, generating explanations,
and proposing candidate transformations. They are never treated as authorities over
financial truth.

All semantic validity is enforced by deterministic compilers and invariant checks.
An agent-orchestrated control plane coordinates intent capture, validation, compilation,
simulation, and approval, ensuring that machine learning components cannot violate core
accounting guarantees.

\section{System Organization and Governance}

The system is organized as a layered, hermetic workspace in which the accounting kernel,
intermediate representations, and compilers evolve together under strong dependency
constraints. Governance services, authoring interfaces, and runtime deployment are
decoupled from this core via signed, content-addressed artifacts.

This structure enables rapid evolution while preserving correctness guarantees,
auditability, and long-term reproducibility.

\section{Contributions and Implications}

This architecture synthesizes principles from programming languages, database systems,
and computer architecture into a unified substrate for institutional accounting. By
treating accounting as deterministic conservation over immutable inputs, the system
achieves high assurance and high performance simultaneously, without sacrificing
extensibility or transparency.

\end{document}